\documentclass[11pt,a4paper]{article}
\usepackage[T1]{fontenc}
\usepackage[utf8]{inputenc}
\usepackage{amsmath,amssymb,amsthm}
\usepackage[margin=1in]{geometry}
\usepackage[colorlinks=true,linkcolor=blue,urlcolor=blue]{hyperref}
\theoremstyle{plain}
\newtheorem{theorem}{Theorem}
\newtheorem{lemma}[theorem]{Lemma}
\newtheorem{proposition}[theorem]{Proposition}
\newtheorem{corollary}[theorem]{Corollary}
\theoremstyle{definition}
\newtheorem{remark}[theorem]{Remark}
% A quoted conjecture can be a single hundred-term formula whose terms are thirty-digit
% numbers. TeX cannot hyphenate those, so without a little stretch every such quote runs off
% the page; the linked-a-file papers are all of that shape.
\setlength{\emergencystretch}{3em}

\title{The empirical closed form for OEIS A219911, proved}
\author{Adrian Perez Fontelles\\ \small Independent researcher}
\date{1 September 2026}
\begin{document}
\maketitle

\begin{abstract}
OEIS A219911 carries an empirical closed form for $a(n)$ --- not a recurrence relating
consecutive terms but an explicit formula. It is true, and it is decidable rather than
empirical. The entry counts arrays subject to a condition on a bounded window of consecutive lines, so the lines of an array are the
vertices of a finite digraph, an array is a walk in it, and $a(n)$ is a walk count on
$S=483$ vertices. Any function of the form $\sum_j p_j(n)\lambda_j^{\,n}$ satisfies the
monic linear recurrence whose characteristic polynomial is
$\prod_j(x-\lambda_j)^{\deg p_j+1}$; here that polynomial is $q(x)=(x-1)^{30}$, of degree
$30$. Two things are then checked exactly: that the walk count satisfies the same
recurrence from some index on, which the Cayley--Hamilton theorem reduces to a finite
computation, and that the two sequences agree at $30$ consecutive indices past that
point. A monic recurrence determines every later term from its predecessors, so agreement
there is agreement for good.
\end{abstract}

\noindent\small 2020 Mathematics Subject Classification. 05A15, 11B37, 15A18.\normalsize

\section{The conjecture}

OEIS A219911 is ``Number of nX4 arrays of the minimum value of corresponding elements and their horizontal, vertical, diagonal or antidiagonal neighbors in a random, but sorted with lexicographically nondecreasing rows and nonincreasing columns, 0..2 nX4 array.''. It has offset $1$ and begins
\[
10, 10, 64, 234, 819, 2563, 7877, 24318,\ \dots
\]
The entry states, as a conjecture contributed by its author and never marked settled:
\begin{quote}\raggedright\small
Empirical: a(n) = (1/8841761993739701954543616000000)*n\^{}29 - (1/21777738900836704321536000000)*n\^{}28 + (61/6222211114524772663296000000)*n\^{}27 - (43/31425308659216023552000000)*n\^{}26 + (17/125132450769728962560000)*n\^{}25 - (28109/2863608007999566643200000)*n\^{}24 + (790000019/1563529972367763387187200000)*n\^{}23 - (7238713/441425740363569561600000)*n\^{}22 + (110073307/882851480727139123200000)*n\^{}21 + (4489349/230992014842265600000)*n\^{}20 - (543475622663/462446013714215731200000)*n\^{}19 + (2560006651583/85187423578934476800000)*n\^{}18 + (2310801616411997/46512333274098224332800000)*n\^{}17 - (589559743689019/17766361067264409600000)*n\^{}16 + (496003326175726703/390859943479817011200000)*n\^{}15 - (209411840851547611/9306189130471833600000)*n\^{}14 + (8651366663182657/468062766917222400000)*n\^{}13 + (88016360529632233/9308066387558400000)*n\^{}12 - (26031638075762884878329/105340233495851827200000)*n\^{}11 + (120380686054081150949947/41383663159084646400000)*n\^{}10 - (366096865860529707276691/133888321985273856000000)*n\^{}9 - (27222349343899273729289/57898389136896000000)*n\^{}8 + (4251094114198631270606723/540988073497872000000)*n\^{}7 - (148882500240267881827577429/2448107645727744000000)*n\^{}6 + (17631394572311382143692779449/122527787668673587200000)*n\^{}5 + (465421984465716632128098689/291732827782556160000)*n\^{}4 - (241713358053685840278268817/14586641389127808000)*n\^{}3 + (706593057831502942631/10962783631800)*n\^{}2 - (8901027387557259947/93163582512)*n - 5119093 for n>18
\end{quote}
As of the ``Last modified'' line on the live entry (Oct 03 2025 18:24:13, revision 8) this is still
recorded as empirical, and nothing on the entry records it as proved. Write
\begin{multline*}
f(n)\;=\;\frac{n^{29}}{8841761993739701954543616000000} - \frac{n^{28}}{21777738900836704321536000000} \\
+ \frac{61 n^{27}}{6222211114524772663296000000} - \frac{43 n^{26}}{31425308659216023552000000} \\
+ \frac{17 n^{25}}{125132450769728962560000} - \frac{28109 n^{24}}{2863608007999566643200000} \\
+ \frac{790000019 n^{23}}{1563529972367763387187200000} - \frac{7238713 n^{22}}{441425740363569561600000} \\
+ \frac{110073307 n^{21}}{882851480727139123200000} + \frac{4489349 n^{20}}{230992014842265600000} \\
- \frac{543475622663 n^{19}}{462446013714215731200000} + \frac{2560006651583 n^{18}}{85187423578934476800000} \\
+ \frac{2310801616411997 n^{17}}{46512333274098224332800000} - \frac{589559743689019 n^{16}}{17766361067264409600000} \\
+ \frac{496003326175726703 n^{15}}{390859943479817011200000} - \frac{209411840851547611 n^{14}}{9306189130471833600000} \\
+ \frac{8651366663182657 n^{13}}{468062766917222400000} + \frac{88016360529632233 n^{12}}{9308066387558400000} \\
- \frac{26031638075762884878329 n^{11}}{105340233495851827200000} + \frac{120380686054081150949947 n^{10}}{41383663159084646400000} \\
- \frac{366096865860529707276691 n^{9}}{133888321985273856000000} - \frac{27222349343899273729289 n^{8}}{57898389136896000000} \\
+ \frac{4251094114198631270606723 n^{7}}{540988073497872000000} - \frac{148882500240267881827577429 n^{6}}{2448107645727744000000} \\
+ \frac{17631394572311382143692779449 n^{5}}{122527787668673587200000} + \frac{465421984465716632128098689 n^{4}}{291732827782556160000} \\
- \frac{241713358053685840278268817 n^{3}}{14586641389127808000} + \frac{706593057831502942631 n^{2}}{10962783631800} \\
- \frac{8901027387557259947 n}{93163582512} - 5119093.
\end{multline*}

\section{The count is a walk count}

The entry's defining condition is local: it is a condition on a bounded window of consecutive lines. Take as vertices the
admissible configurations of such a window and put an edge from $u$ to $v$ when $v$ may
follow $u$, which the condition decides by inspection. An admissible array is then exactly a
walk, so with $M$ the adjacency matrix of that digraph on $S=483$ vertices, and $u,w$ the
vectors recording which windows may start and end an array,
\[
a(n)\;=\;u^{\!\top}M^{\,g(n)}w
\]
for the linear function $g$ that the entry's shape supplies. In particular $a$ is
$C$-finite: it satisfies some monic integer linear recurrence, though not necessarily a short
one.

\section{A closed form is a recurrence}

\begin{lemma}\label{lem:ann}
Let $f(m)=\sum_{j}p_j(m)\lambda_j^{\,m}$ with the $\lambda_j$ distinct and the $p_j$
polynomials, and put $q(x)=\prod_j(x-\lambda_j)^{\deg p_j+1}$. Then $q$ applied to the shift
operator annihilates $f$: writing $q(x)=x^{r}-\sum_{i=1}^{r}c_ix^{r-i}$,
\[
f(m)\;=\;\sum_{i=1}^{r}c_i\,f(m-i)\qquad\text{for every }m.
\]
\end{lemma}

\begin{proof}
The shift operator $E$ acts on $m\mapsto\lambda^{m}p(m)$ as
$(E-\lambda)\bigl(\lambda^{m}p(m)\bigr)=\lambda^{m+1}\bigl(p(m+1)-p(m)\bigr)$, and
$p(m+1)-p(m)$ has degree one less than $p$. So $(E-\lambda)^{\deg p+1}$ kills
$\lambda^{m}p(m)$, and the factors of $q$ commute.
\end{proof}

Here $q(x)=(x-1)^{30}$, of degree $30$; the closed form is a polynomial in $n$, so this is $(x-1)^{30}$, and the recurrence it encodes is
\begin{equation}\label{eq:rec}
a(n)\;=\;30\,a(n-1) - 435\,a(n-2) + 4060\,a(n-3) - 27405\,a(n-4) + 142506\,a(n-5) - 593775\,a(n-6) + 2035800\,a(n-7) - 5852925\,a(n-8) + 14307150\,a(n-9) - 30045015\,a(n-10) + 54627300\,a(n-11) - 86493225\,a(n-12) + 119759850\,a(n-13) - 145422675\,a(n-14) + 155117520\,a(n-15) - 145422675\,a(n-16) + 119759850\,a(n-17) - 86493225\,a(n-18) + 54627300\,a(n-19) - 30045015\,a(n-20) + 14307150\,a(n-21) - 5852925\,a(n-22) + 2035800\,a(n-23) - 593775\,a(n-24) + 142506\,a(n-25) - 27405\,a(n-26) + 4060\,a(n-27) - 435\,a(n-28) + 30\,a(n-29) - a(n-30).
\end{equation}

\section{The proof}

\begin{lemma}\label{lem:crit}
With $q$ as above, $a$ satisfies \eqref{eq:rec} for all large $n$ if and only if
$u^{\!\top}M^{\,j}q(M)w=0$ for every $j\ge0$, and it is enough to check $j<S$.
\end{lemma}

\begin{proof}
Substituting the walk expression, the residual $a(n)-\sum_ic_ia(n-i)$ equals
$u^{\!\top}M^{\,g(n)-r}q(M)w$ up to the fixed linear reparametrisation $g$. For the bound,
$M^{S}$ is an integer combination of $I,M,\dots,M^{S-1}$ by Cayley--Hamilton, so the values
$u^{\!\top}M^{j}q(M)w$ satisfy the monic recurrence given by the characteristic polynomial of
$M$; $S$ consecutive zeros therefore force all later ones, and the last nonzero value pins the
threshold exactly.
\end{proof}

\begin{theorem}
$a(n)=f(n)$ for every $n\ge19$.
\end{theorem}

\begin{proof}
The vector $w'=q(M)w$ was formed by $30$ matrix--vector products in exact integer
arithmetic, and $u^{\!\top}M^{j}w'$ evaluated for $j=0,1,\dots$, again exactly. Those integers
vanish from the point corresponding to $n=49$ onwards, and the run was continued until the
working vector was identically zero, which settles every larger $j$ at once. So $a$ satisfies
\eqref{eq:rec} for every $n>48$. By Lemma~\ref{lem:ann} $f$ satisfies \eqref{eq:rec}
identically. Both were then evaluated in exact arithmetic at the $30$ consecutive indices
$n=49,\dots,78$ and agree at each. A monic recurrence of order $30$
determines $a(m)$ from $a(m-1),\dots,a(m-30)$, and for every $m\ge79$ both
$m>48$ and $m-30\ge49$ hold, so induction on $m$ gives $a(m)=f(m)$ for every
$m\ge49$.
Below $n=49$ the recurrence argument gives nothing, since \eqref{eq:rec} is only
established for $a$ from $n=49$ on. The remaining indices $n=19,\dots,48$ were
therefore checked one at a time, directly against the walk count in exact integer arithmetic,
and agree.
\end{proof}

The range is tight in the sense that was checked: at $n=18$ the closed form and the sequence differ, so no larger range is available.

\section{Verification}

Four checks, each independent of the algebra above.

First, the walk model reproduces every one of the $25$ terms the entry publishes,
exactly, in integer arithmetic. This is what ties the model to the entry: the model is built
by reading the entry's English, and a misreading gives a different digraph and different
counts.

Second, the closed form was evaluated in exact rational arithmetic --- not floating point ---
at every published term from $n=19$ on, and agrees with the entry's own DATA at each.

Third, the annihilation test of Lemma~\ref{lem:crit} was carried out exactly, and the model
and the closed form were then compared directly at sixty consecutive indices beyond
$n=19$, far past where the proof already settles the matter.

Fourth, the closed form was deliberately perturbed --- by a constant and by a term linear in
$n$ --- and each perturbed form was rejected by the same comparison, so the test is not one
that anything would pass.

\begin{thebibliography}{9}
\bibitem{oeis} The OEIS Foundation, \emph{The On-Line Encyclopedia of Integer
Sequences}, \url{https://oeis.org}, 2026. Sequence A219911.
\bibitem{stanley} R.~P.~Stanley, \emph{Enumerative Combinatorics}, Volume 1, 2nd ed.,
Cambridge University Press, 2011. (Transfer-matrix method, Section 4.7; $C$-finite sequences,
Section 4.1.)
\bibitem{kp} M.~Kauers and P.~Paule, \emph{The Concrete Tetrahedron}, Springer, 2011.
(Closed forms and linear recurrences with constant coefficients.)
\bibitem{horn} R.~A.~Horn and C.~R.~Johnson, \emph{Matrix Analysis}, 2nd ed., Cambridge
University Press, 2013.
\end{thebibliography}

\end{document}
